\documentclass[11pt,a4paper]{article}
\usepackage[margin=1in]{geometry}
\usepackage[T1]{fontenc}
\usepackage{lmodern}
\usepackage{microtype}
\usepackage{hyperref}
\usepackage{array}
\usepackage{longtable}
\hypersetup{
  colorlinks=true,
  linkcolor=black,
  urlcolor=black,
  pdftitle={OC Core 1.3.2 Journal Core},
  pdfauthor={Alexander Yashin},
  pdfsubject={Ontology of Continua Core 1.3.2 journal-core release artifact},
  pdfkeywords={Ontology of Continua, OC Core 1.3.2, release, journal core}
}
\title{\textbf{OC Core 1.3.2 Journal Core}\\
\large Release-Bound Scientific Kernel for Collapse, Residue, Rebirth, and Evidence-Scoped Continuity}
\author{Alexander Yashin\\Independent Researcher, Leipzig/Halle, Germany\\ORCID 0009-0008-6166-0914}
\date{Version v1.3.2; release date: 28.04.2026}
\begin{document}
\maketitle

\begin{abstract}
This journal-core artifact states the compact scientific kernel of OC Core 1.3.2. It is not a replacement for the master monograph; it is the release-bound article spine that a reviewer can read first. The central result is a source-audited grammar for collapse, residue, and rebirth of continua, together with explicit evidence boundaries for theorem-native, operational, simulation, and frontier claims. OC Core 1.3.2 includes the terminal release-blocker ledger, K0--K12 projection state, DRT strict-salvage state, Tsukrov repair, Parfitian Cerberus pass, LRGEF release state, and night-delta integration as bounded public-safe release artifacts.
\end{abstract}

\section*{Release Identity}
OC Core 1.3.2 is the current stable no-send baseline of the Core 1.3 line. It is prepared for owner review and package audit, not for automatic public publication. Public release, DOI minting, Zenodo upload, GitHub release creation, and outbound announcements remain locked until explicit owner approval and publication-gate review.

\section*{Scientific Kernel}
The release kernel has four load-bearing notions:
\begin{itemize}
  \item admissible realization: a state or fiber region that keeps a continuum live under its declared threshold and embedding regime;
  \item collapse/death: the loss of every admissible realization under that regime;
  \item residue: post-collapse structure that may remain without preserving live identity;
  \item rebirth: a later live continuum grounded in residue but not identical with the original continuum.
\end{itemize}
The defended classification claim is therefore bounded: within the audited OC Core kernel, when a continuum loses every admissible realization, the original continuum is dead; surviving structure is residue, and later residue-grounded life is rebirth rather than persistence of the original continuum.

\section*{Evidence Boundary}
OC Core 1.3.2 separates claim support into release-visible classes. The master monograph carries theorem-roadmap references and proof-machinery anchors; the release packet carries evidence ledgers, replay manifests, domain projection state, and reviewer-navigation materials. The package does not claim unbounded predictive authority, full replacement of domain science, or completion of every possible extension theorem.

\section*{Release-Integrated Science}
\begin{longtable}{>{\raggedright\arraybackslash}p{0.28\linewidth}>{\raggedright\arraybackslash}p{0.64\linewidth}}
\textbf{Surface} & \textbf{Role in 1.3.2}\\
\hline
Master monograph & Canonical theory, definitions, theorem roadmap, proof machinery, and evidence boundaries.\\
Science blocker closure ledger & Release terminality ledger for the eighty-two tracked release blockers; it records proof/evidence, replay, and owner-gate terminal routes without treating mere administrative notes as science.\\
K0--K12 projection state & Domain projection/replay status across mathematics, physics, chemistry, biology, systems/civilizational projection, and metaontology.\\
DRT strict salvage & Strict-salvage status: DRT-derived material is not promoted into Core authority without reproducibility and Core-gate reconciliation.\\
Tsukrov repair & Evidence-first response surface for the Stanislav Tsukrov review packet, using bounded public summaries and release-safe references.\\
Parfitian Cerberus & Ethical and release-governance audit layer for no-send release readiness.\\
LRGEF & Release governance, manifest, provenance, ZIP integrity, version consistency, and no-send publication lock.\\
\end{longtable}

\section*{What This Article Does Not Claim}
This article imports the public-safe formal and benchmark elements prepared for the OC 1.4 line when they already have proof, replay, or gate support: observation architecture, event/fact contract, support classes, quantum/no-signalling boundary, IMN baseline treatment, and the deterministic ten-task benchmark suite. It does not include private IP material, private reviewer-source material, private model transcripts, account credentials, local filesystem paths, or draft-only internal research. It does not use unsupported theorem labels or unsupported prediction language. Quantitative claims require protocol, code, data or fixture, output hash, tolerance envelope, and falsifier.

\section*{Reviewer Reading Order}
A reviewer should start with this journal core, then read the master monograph front matter and theorem roadmap, then inspect the release dossier, PDF quality report, science blocker closure ledger, K0--K12 projection packet, DRT strict-salvage packet, Tsukrov repair packet, and Parfit/LRGEF gate reports. That reading order keeps the scientific claim, the evidence route, and the publication lock visible at the same time.

\section*{Reader-Specific Navigation}
\begin{longtable}{>{\raggedright\arraybackslash}p{0.25\linewidth}>{\raggedright\arraybackslash}p{0.65\linewidth}}
\textbf{Reader} & \textbf{Primary question}\\
\hline
Owner & Is the package honest enough for public review after the no-send locks are explicitly released?\\
Mathematical reviewer & Are identity, collapse, residue, and rebirth tied to stable definitions and theorem-fate routes?\\
Domain reviewer & Are the domain projections bounded by replay state rather than treated as domain replacement?\\
Reproducibility reviewer & Can every package claim be traced to source, manifest, hash, data, or gate output?\\
Publication reviewer & Are all external actions still locked until owner and channel approval?\\
\end{longtable}

\section*{Proven, Replayed, Future}
\begin{longtable}{>{\raggedright\arraybackslash}p{0.28\linewidth}>{\raggedright\arraybackslash}p{0.58\linewidth}}
\textbf{Class} & \textbf{Release treatment}\\
\hline
Proven or evidence-closed & The Core kernel, theorem-fate rows, source-audited support classes, and terminal blocker ledger.\\
Replayed & K0--K12 domain projection routes, benchmark/replay manifests, simulation-result records, and DRT workflow/salvage checks.\\
Outside 1.3.2 & IP-sensitive invention review, unrelated-domain claims, broader empirical protocols, and owner-gated publication actions.\\
\end{longtable}

\section*{Minimal Scientific Claim}
The minimal scientific claim of this release is not that OC has finished every possible theorem. It is that the Core 1.3.2 package now exposes a coherent, evidence-scoped ontology of continua with release-visible boundaries. A reviewer can inspect where a claim is formal, where it is replayed, where it is structural, and where a later release must do more work.

\section*{Publication Boundary}
The journal core is publication-ready only as part of the owner-reviewed release package. It does not authorize publication by itself, and it does not widen any claim beyond the release manifest, claim ledger, support map, and no-send governance state.

\clearpage
\section*{Release Inspection Appendix A: Claim Support}
\begin{longtable}{>{\raggedright\arraybackslash}p{0.24\linewidth}>{\raggedright\arraybackslash}p{0.66\linewidth}}
\textbf{Claim type} & \textbf{Required support in 1.3.2}\\
\hline
Identity and continuity & Definition route in the master monograph, source-audited terminology, and theorem-fate row when stated formally.\\
Collapse/death & Explicit admissible-realization criterion and support in the Core kernel.\\
Residue & Bounded post-collapse classification, with no claim that residue preserves live identity.\\
Rebirth & Later live structure grounded in residue, explicitly not identical with the collapsed original.\\
Domain projection & K/domain packet, replay status, source-native inclusion verdict, and structural-only boundary where quantitative replay is absent.\\
Publication readiness & LRGEF, Parfitian Cerberus, manifest/checksum integrity, and owner no-send lock.\\
\end{longtable}

\clearpage
\section*{Release Inspection Appendix B: Reviewer Questions}
\begin{enumerate}
  \item Does the release define its central terms before using them as claims?
  \item Does each strong theorem, prediction, or promotion word have a support-map row?
  \item Does each domain extension remain bounded by its K/domain projection state?
  \item Does the package keep private review material and model transcripts outside the public bundle?
  \item Does the publication package remain no-send until owner/channel unlock?
\end{enumerate}
Positive answers to these questions do not make OC final. They make this release inspectable, citable after approval, and scientifically bounded.

\clearpage
\section*{Release Inspection Appendix C: What Would Be Misleading}
It would be misleading to market this release as an unbounded predictive engine or as authority over unrelated problem classes. The correct public posture is strong and precise: OC Core 1.3.2 includes a global verdict-invariant minimality theorem for OC-compatible representations, deterministic benchmark support for the admitted forward layer, visible evidence routes, a terminal blocker ledger, and explicit owner/publication locks.

\end{document}
